\subsection{Poincar\'e Duality}

Let $S\subseteq M$ be an oriented $k$-submanifold of an oriented $n$-manifold $M$, and further suppose $S$ is closed in $M$.
For a closed $k$-form $\omega\in\Omega_c^k(M)$, we have that $\supp(\omega|_S)$ is closed in $S$ and thus in $M$, and it is a closed subset of $\supp\omega\cap S$ which is then compact.
Thus $\omega|_S$ is compactly supported in $S$ and so can be integrated.
Thus we have defined a linear functional
$$ H^k_c(M) \to \bR,\qquad \omega \mapsto \int_S\omega|_S . $$
Because $(H^k_c(M))^*\cong H^{n-k}(M)$ by Poincar\'e duality, this functional is represented by some $\eta_S\in H^{n-k}(M)$.
That is to say, for $\omega\in H^k_c(M)$,
$$ \int_S\omega|_S = \int_M\omega\wedge\eta_S . $$
This $\eta_S$ is called the {\emphcolor Poincar\'e dual} of $S$.

Now if $S$ is compact, then it is closed and as before we have $\eta_S$ a {\it closed\/} Poincar\'e dual.
Because $S$ is compact we can integrate any $k$-form on $M$.
Thus we have a functional
$$ H^k(M) \to \bR,\qquad \omega\mapsto\int_S\omega|_S . $$
If $M$ is of finite type, then $H^k(M)^*\cong H^{n-k}_c(M)$, and so there is a $\eta_S'\in H^{n-k}_c(M)$ such that
$$ \int_S\omega|_S = \int_M\omega\wedge\eta_S'\qquad\text{for all $\omega\in H^k(M)$} . $$
Note that if this holds for all $\omega$, it holds for all compactly supported $\omega$ and thus $\eta'_S$ is also the closed Poincar\'e dual.
So the natural map $H^\bul_c(M)\to H^\bul(M)$ sends the compact Poincar\'e dual to the closed Poincar\'e dual.
But their cohomology classes may be very different.

For instance in $\bR^n$, the closed Poincar\'e dual of any point $p$ is trivial since $H^n(\bR^n)=0$.
But the compact Poincar\'e dual has integral $1$.

We will focus on the closed Poincar\'e dual:
$$ \int_S\omega|_S = \int_M\omega\wedge\eta_S \qquad\text{for $\omega\in H^k_c(M)$} . $$
Recall that if $S\subseteq M$ is a submanifold, then $S$ has a {\emphcolor tubular neighborhood}.
That is, there is a vector bundle $\pi\colon N\to S$ and a smooth map $j\colon N\to M$ such that $j(p,0)=p$ and $j$ maps an open neighborhood of the zero section of $N$ diffeomorphically
onto an open neighborhood of $S$.
This $N$ can be taken to be the normal bundle of $S$, that is $N_S=TM|_S/TS$.
And furthermore the tubular neighborhood may be taken to be all of $N_S$.

Note that if $A,B$ are oriented vector bundles with oriented trivializations $(U_i,\phi_i)$ and $(U_i,\psi_i)$, then $A\oplus B$ is oriented with $(U_i,\phi_i\oplus\psi_i)$.
If $T\subseteq N_S$ is a tubular neighborhood of $S$, then as $N_S$ is the quotient of the oriented $TM|_S$ by the oriented $TS$ (we assume $S$ is orientable), it too is orientable.
Furthermore we have
$$ N_S\oplus TS = TM|_S , $$
so we oriented $N_S$ in such a way to make the orientation on $TM|_S$ the direct sum orientation.

Consider the Thom isomorphism on $N_S\to S$
So we have a composition
$$ H^\bul(S) \xtos{\bul\wedge\Phi}[\sim] H^{\bul+n-k}_\cv(N_S) \xtos{j_*} H^{\bul+n-k}(M) . $$
Here $j_*$ extends a form on $N_S$ by zero on all of $M$.
We claim that
$$ \eta_S = j_*(\Phi) \in H^{n-k}(M) . $$
Now let $\omega\in H^k_c(M)$, let $i\colon S\to N_S$ the zero section (viewed as inclusion into the tubular neighborhood), and $\pi\colon N_S\to S$ the projection.
Since $\pi$ is a deformation retract, so that $\pi,i$ are homotopy inverses and thus $\pi^*,i^*$ are inverses in cohomology.
Thus $\pi^*i^*\omega$ is cohomologous with $\omega$ on $T$: they differ by an exact form $\omega=\pi^*i^*\omega+\d\tau$.

Now notice

\medskip
\centerline{\openup1\jot\vbox{\halign{\hfil$\displaystyle#$&$\displaystyle{}=#$\hfil\quad&#\hfil\cr
    \int_M\omega\wedge j_*\Phi&\int_N\omega\wedge\Phi&because $j_*\Phi$ has support in $N$\cr
        &\int_N(\pi^*i^*\omega+\d\tau)\wedge\Phi\cr
        &\int_N\pi^*i^*\omega\wedge\Phi&because $(\d\tau)\wedge\Phi=\d(\tau\wedge\Phi)$ has zero integral by Stokes'\cr
        &\int_Si^*\omega\wedge\pi_*\Phi&by \refmath[proposition]{projformulas} (2)\cr
        &\int_S\omega|_S&since $\pi_*\Phi=1$ by definition.\cr
}}}
\medskip

Thus $j_*\Phi$ satisfies the defining property of the Poincar\'e dual, so $\eta_S=j_*\Phi$.
Note that if $S$ is compact then this means $\eta_S$ has compact support.

Now suppose $E\to M$ is an oriented vector bundle.
Then $M$ embeds into $E$ via the zero section, so it is a submanifold.
There is an exact sequence
$$ 0 \to TM \to TE|_M \to E \to 0 , $$
meaning the normal bundle of $M$ in $E$ is $E$ itself.
Thus as we have shown, the Poincar\'e dual of $M$ in $E$ is $E$'s Thom class.

We summarize.

\bprop

    \benum
        \item The Poincar\'e dual of a closed (as a subspace) oriented submanifold $S\subseteq M$ can be identified with the Thom class of $S$'s normal bundle.
        \item The Thom class of a vector bundle $E\to M$ and the Poincar\'e dual of the zero section of $E$ may be identified.
    \eenum

\eprop

If $R,S\subseteq M$ are submanifolds which intersect transversally, then their normal bundles satisfy
$$ N_{R\cap S} = N_R\oplus N_S , $$
and thus their Thom classes satisfy
$$ \Phi(N_{R\cap S}) = \Phi(N_R\oplus N_S) = \Phi(N_R)\wedge\Phi(N_S) $$
and thus the Poincar\'e duals satisfy
$$ \eta_{R\cap S} = \eta_R\wedge\eta_S . $$

